\section{模型假设}

结合附件覆盖范围与题设成本，我们作出以下假设：

\begin{itemize}[itemindent=2em]
\item \textbf{假设1：} 在附件 A 的训练配方范围内，17 个领域占比的线性组合能够近似描述验证损失；换到其他参数规模时，允许损失的整体尺度发生变化。
\item \textbf{假设2：} B6、B7 的质量随评分提高而提高，质量不足带来的损失可作为单独一项加到经典标度律中。其他来源须先核对评分含义和方向，再确定用法。
\item \textbf{假设3：} 按题设将总算力分为基础训练、质量提升和长文本注意力三项；质量费用采用附录 B 的函数，上下文采用 C7 的给定长度，不参与自由优化。
\item \textbf{假设4：} 当前窗口的能力前沿可由 logistic 函数近似，规模与得分的条件关联可由含年份项的合并横截面回归描述；沿用这些关系的未来结果视为条件情景。
\end{itemize}

线性近似用于比较领域系数，核岭用于同尺度预测和训练凸包内的配方选择。相同配方迁移到其他参数规模时，仍需检验其损失水平；模型选择与支持范围约束均不能替代新配方的训练实验。

加性质量项便于比较各类投入的边际作用，其与规模可分离的形式属于建模假设。附件 A 的综合评分与附件 B 的实验质量变量并无配对标定，问题三暂按两者数值相同传递基线，并另用单调映射作敏感性对照。评分方向一致只能支持排序，不能证明两种尺度相同。

预算优化沿用题设费用体系。参数量和数据量共同决定基础训练开销，质量费用随处理量累计，上下文长度则作为给定情景进入注意力开销。由此得到的配置是这些条件共同作用的结果，更换费用关系时应同时观察配置与损失的变化。

C1 的不同年份记录没有形成同一模型的跨年追踪，因而回归系数描述的是模型之间的条件差异。前沿分解将规模关联项以外的增量记为剩余项，其中也含样本构成和遗漏因素。前沿曲线则描述高分记录的累积变化；我们保留其观测窗口，并通过窗口删减和时间留出检查外推对平台段的依赖。
